\section*{第 \ref{ch:g4:heart-and-blood} 章——心脏与血液}

\begin{solution}{exo:g4:heart-and-blood:1}
\omterm{def:g4:journey-of-food:digestion}{营养物质}在\omterm{def:g4:journey-of-food:tract}{小肠}中被吸收；氧气
在肺部被吸收。
\end{solution}

\begin{solution}{exo:g4:heart-and-blood:2}
二氧化碳——在肺部被释放，
并随\omterm{def:g4:breathing:movements}{呼气}排出体外。
\end{solution}

\begin{solution}{exo:g4:heart-and-blood:3}
它是一块空心\omterm{def:g3:bones-and-muscles:muscle}{肌肉}，大小约与主人的拳头相当，
位于胸腔正中，紧贴\omterm{def:g3:bones-and-muscles:skeleton}{肋骨}之后。
\end{solution}

\begin{solution}{exo:g4:heart-and-blood:4}
右侧\omterm{def:g4:heart-and-blood:heart}{心脏}将\omterm{prop:g4:journey-of-food:absorb}{血液}泵送至肺部以获取氧气；
左侧则将富含氧气的\omterm{prop:g4:journey-of-food:absorb}{血液}泵送至全身。
\end{solution}

\begin{solution}{exo:g4:heart-and-blood:5}
每次\omterm{def:g4:heart-and-blood:heart}{心跳}在\omterm{def:g4:heart-and-blood:vessels}{血管}中推动一股小波动，
可在腕部或颈部摸到。数它即是在数
每分钟的\omterm{def:g4:heart-and-blood:heart}{心跳}次数。
\end{solution}

\begin{solution}{exo:g4:heart-and-blood:6}
拇指本身有明显的小\omterm{def:g4:heart-and-blood:heart}{搏动}，
因此可能会数成拇指的\omterm{prop:g4:heart-and-blood:pulse}{脉搏}，而非所测的腕部\omterm{prop:g4:heart-and-blood:pulse}{脉搏}。
\end{solution}

\begin{solution}{exo:g4:heart-and-blood:7}
短跑时腿部快速消耗\omterm{def:g4:journey-of-food:digestion}{营养物质}和氧气，
必须迅速补充；\omterm{def:g4:heart-and-blood:heart}{心脏}加快并加强\omterm{def:g4:heart-and-blood:heart}{搏动}，
使\omterm{prop:g4:journey-of-food:absorb}{血液}更频繁地往返输送。
\end{solution}

\begin{solution}{exo:g4:heart-and-blood:8}
答案不唯一。静息时通常约为 $80$--$100$，
开合跳后约为 $120$--$150$——输送系统正在发挥作用。
\end{solution}

\begin{solution}{exo:g4:heart-and-blood:9}
刚离开肺部时，\omterm{prop:g4:journey-of-food:absorb}{血液}刚装载了氧气——颜色鲜红。
在\omterm{def:g3:bones-and-muscles:muscle}{肌肉}中，它交出氧气，并带走
二氧化碳——颜色较暗。颜色随所运物质而定。
\end{solution}

\begin{solution}{exo:g4:heart-and-blood:10}
\omterm{def:g4:journey-of-food:tract}{小肠}将\omterm{def:g4:journey-of-food:digestion}{营养物质}送入\omterm{prop:g4:journey-of-food:absorb}{血液}；肺为\omterm{prop:g4:journey-of-food:absorb}{血液}补充氧气、
并排出二氧化碳；\omterm{def:g4:heart-and-blood:heart}{心脏}则是泵，
使\omterm{prop:g4:journey-of-food:absorb}{血液}在这些器官之间流动，并输送到各个
需要物质的部位：\omterm{def:g3:bones-and-muscles:muscle}{肌肉}、大脑、\omterm{def:g3:bones-and-muscles:skeleton}{骨骼}和\omterm{def:g1:the-five-senses:sense}{皮肤}。
\end{solution}

\begin{solution}{exo:g4:heart-and-blood:11}
经过锻炼的\omterm{def:g4:heart-and-blood:heart}{心脏}是更强壮的\omterm{def:g3:bones-and-muscles:muscle}{肌肉}：每次\omterm{def:g4:heart-and-blood:heart}{搏动}输送的\omterm{prop:g4:journey-of-food:absorb}{血液}更多。
输送同样的血量所需的每分钟\omterm{def:g4:heart-and-blood:heart}{搏动}次数便减少了——
因此骑车者平静时的 $50$ 次\omterm{def:g4:heart-and-blood:heart}{搏动}，相当于未经训练者的 $75$ 次，
余力留作剧烈运动之用。
\end{solution}
